\par\Needspace{9\baselineskip}\bigskip
\pdfbookmark[3]{Darstellung, [24]}{pallas-II-article24}
\begin{center}[24].\end{center}
\begingroup
\makeatletter
\def\@makefnmark{\hbox{\normalfont\@thefnmark}}
\makeatother
\noindent\hspace*{\parindent}Bei einem rationalen Verhältnis von \(n\text{ und }n'\) enthält der Wert von \(n+\delta n\) einen der Zeit proportionalen Teil oder eine säkulare Änderung, und die mittlere Bewegung \(\int n\mathrm{d}t\) enthält ein dem Quadrat der Zeit proportionales Glied. Für die gegenwärtigen Bedürfnisse der Astronomie kann man es dabei belassen. Wir werden jedoch sehen, dass die Betrachtung der von höheren Potenzen der Massen abhängigen Glieder dieser dem Anschein nach säkularen Gleichung eine \textit{periodische Form} zurückgibt. Dies ist umso bemerkenswerter, als uns die Theorie der Bewegung der Pallas diesen Fall tatsächlich liefern wird\HistoricalNoteEight{Siehe S. 557---559.] }.
\par\endgroup
